%% ---------------------------------------------------------------------
\section{Introduction}\label{sec:intro}

The firing squad synchronization problem (FSSP), posed by Myhill in 1957
and popularised by Moore~\cite{Moore64}, asks for a finite automaton that,
placed in every cell of a line of arbitrary length $n$ and started with all
cells quiescent except a \emph{general} at one end, drives all cells into a
common \emph{firing} state for the first time at the same instant. Goto
\cite{Goto62} and Waksman \cite{Waksman66} showed that this is possible in
time $2n-2$, which is optimal \cite{Moore64}. Since then the minimum number
of states of a \emph{minimal-time} solution has been a benchmark question
in cellular automata theory \cite{Mazoyer86,Umeo09}: Waksman used $16$
states, Balzer $8$ \cite{Balzer67}, Gerken $7$ \cite{Gerken87} and Mazoyer
$6$ \cite{Mazoyer87}. On the other side, Balzer \cite{Balzer67} reported,
and Sanders \cite{Sanders94} confirmed with a corrected exhaustive search,
that no $4$-state minimal-time solution exists. Whether five states suffice
has remained open since Mazoyer's construction; it is regarded as the
central open question on state complexity of the FSSP
\cite{Umeo09,Correa17}. Many six-state minimal-time solutions are known
today \cite{Clergue18,NguyenMaignan20}, and five-state solutions are known
for restricted families of lengths or for non-minimal time
\cite{UmeoYanagihara07,Yunes08,UmeoKamikawaYunes09}, but no five-state
minimal-time solution has been found, nor has its existence been refuted.

\paragraph{This paper.}
We develop tools that apply to minimal-time solutions with any number of
states and use them to attack the five-state question.
\begin{enumerate}[label=(\roman*)]
\item \emph{Agreement and reflection.} We make explicit the elementary but
  useful fact that the line of length $n$ coincides with the half-line below
  the anti-diagonal $t+i=2n-2$, so that all lines are encoded by a single
  half-line diagram $C_\infty$ plus one ``reflected triangle'' $R_n$ per
  length, and that $R_n$ is a function of two anti-diagonals of $C_\infty$
  only (\cref{sec:prelim,sec:pumping}).
\item \emph{A pumping lemma.} We prove that in every minimal-time solution
  the input words of the reflected triangles are pairwise distinct on a cone
  of length about $n/2$ (\cref{lem:pumping}). This is a constraint on the
  half-line alone, valid for \emph{all} pairs of lengths, and it implies a
  \emph{speed-$1/3$ barrier} (\cref{thm:third}): the periodic regime behind
  the front can never overtake a signal of speed $1/3$ emitted by the
  general. This gives an a priori explanation of the speed-$1/3$ signal
  present in all classical constructions.
\item \emph{A certified SAT framework.} We encode ``there is a $k$-state rule
  that synchronizes a given set of lengths in minimal time, together with
  the pumping constraints and further consequences'' into propositional
  logic (\cref{sec:encoding}). The encoding re-derives the non-existence of
  four-state minimal-time solutions in less than a minute on a laptop core,
  and it is validated against Mazoyer's six-state solution extracted from
  the Coq formalization of Duprat \cite{Duprat97}.
\item \emph{Five states.} We report the outcome of the search for five-state
  solutions (\cref{sec:experiments}).
\end{enumerate}

\paragraph{Related work.}
Surveys of the FSSP are \cite{Mazoyer86,Umeo09}. Beyond Balzer and Sanders,
lower bounds on the number of states are known for variants: no three-state
solution and no four-state symmetric minimal-time solution exist for rings
\cite{BerthiaumeEtAl04}; for further restricted classes see
\cite{Correa17}. Computer searches have been used to find new
solutions \cite{Clergue18} and to generate families of solutions by local
simulations \cite{NguyenMaignan20}. SAT solvers have been applied to many
extremal problems in combinatorics with certified unsatisfiability proofs
\cite{HeuleKullmannMarek16,Heule18}; to our knowledge this is the first
SAT-based treatment of the FSSP state-complexity question.
